\chapter{IB-ME方案设计}

\section{两类原语怎样配合}
HIBE允许主密钥为第一层身份发钥，第一层密钥再为下一层身份委派密钥。本作品把接收者$r$放在第一层、发送者$s$放在第二层。因此，持有接收密钥$\mathsf{dk}_r$的用户，可以为期望发送者$s$产生叶钥$\mathsf{dk}_{r,s}$，并尝试解开路径$(r,s)$上的密文。

IBS负责认证发送者。发钥机构为身份$s$生成签名密钥$\mathsf{sk}_s$，验证者以公开参数和身份字符串验证消息签名。组合后，HIBE负责按双方身份打开载荷，IBS负责验证恢复出的载荷确由该发送身份签发。

表~\ref{tab:notations}列出主要符号。两类原语分别维护自己的参数与主秘密，模数和内部表示可以不同。

\begin{table}[H]
\centering\small
\caption{协议主要符号}
\label{tab:notations}
\begin{tabularx}{\textwidth}{@{}p{.23\textwidth}X@{}}
\toprule
\rowcolor{black!8} 符号 & 含义 \\
\midrule
$r,s$；$r',s'$ & 指定的接收者与实际发送者；接收密钥身份与期望发送者 \\
$\mathcal H,\mathcal S$ & 两层匿名HIBE与IBS \\
$\mathsf{hpk},\mathsf{ipk}$ & HIBE和IBS公开参数 \\
$\mathsf{dk}_r,\mathsf{sk}_s$ & 接收者一级密钥与发送者签名密钥 \\
$M,\sigma$ & 应用消息与IBS签名 \\
$z,K$ & 32字节随机种子与派生的对称密钥 \\
$\mathsf{ovk},\mathsf{osk},\tau$ & 一次性验签公钥、签名私钥与签名 \\
$\mathrm{LP}(x)$ & 4字节长度前缀与字节串$x$的拼接 \\
\bottomrule
\end{tabularx}
\end{table}

\section{建立、发钥和基本收发过程}
系统建立时，分别运行HIBE和IBS的建立算法，发布$\mathsf{hpk}$与$\mathsf{ipk}$，由发钥机构保管两份主秘密。发送者发钥调用IBS的身份密钥生成算法；接收者发钥调用HIBE的第一层提取算法。同一身份后续使用同一份已发密钥。

先考虑Wang通用构造的基本收发关系~\cite{wang2024ibme}。发送者对接收身份与消息的规范编码签名，再把签名和消息送入两层HIBE：
\begin{align*}
 \sigma&\leftarrow\mathcal S.\mathsf{Sign}
       (\mathsf{sk}_s,\operatorname{EncMsg}(r,M)),\\
 P&\leftarrow\operatorname{EncPayload}(\sigma,M),\\
 C&\leftarrow\mathcal H.\mathsf{Encrypt}(\mathsf{hpk},(r,s),P).
\end{align*}
接收者使用$\mathsf{dk}_{r'}$为$s'$委派叶钥，解出载荷后，以$s'$验证$\operatorname{EncMsg}(r',M)$的签名。加密路径与验证消息共同绑定双方身份。后面的混合封装保持这一收发关系，并解决载荷长度和传输字节认证问题。

\section{短种子封装应用消息}
IBS签名可能达到数KiB甚至数MiB，直接用短消息HIBE加密整个载荷会超过容量。本作品把HIBE输入固定为32字节随机种子$z$，从$z$派生对称密钥$K$，再使用AES-256-GCM加密签名与应用消息：
\begin{align*}
 C_0&\leftarrow\mathcal H.\mathsf{Encrypt}
       (\mathsf{hpk},(\operatorname{Pad}_P(r),\operatorname{Pad}_C(s)),z),\\
 K&\leftarrow\operatorname{HKDF\text{-}SHA256}(z;\mathsf{ctx}_H,\mathsf D_{\rm kdf}),\\
 C_1&\leftarrow\operatorname{AES\text{-}256\text{-}GCM.Enc}(K,n,P;\mathsf{AAD}).
\end{align*}
这里$n$为新生成的12字节nonce。HKDF和GCM分别采用RFC~5869与NIST SP~800-38D规定的原语~\cite{rfc5869,nist2007gcm}。附加认证数据绑定公开上下文、HIBE胶囊$C_0$、nonce及载荷长度。每条消息重新生成$z$与$n$。

外部身份使用角色标识、长度前缀和零填充编码为定长输入。当前64字节外部身份对应86字节内部身份；接收者与发送者使用不同角色字节，身份哈希摘要不作为公开密文字段。可变长度消息由对称加密承载，HIBE始终处理固定长度种子。

\section{完整密文的字节认证}
对加密数据，系统既要验证发送身份，也要检查传输格式和完整密文字节。为此，每次加密生成一把Lamport型一次性签名密钥~\cite{lamport1979ots}，将一次性公钥的承诺与应用消息共同纳入IBS签名：
\begin{align*}
 d_{\rm ovk}&=H_{\rm commit}(\mathsf{ovk}),\\
 E&=\mathsf D_{\rm bound}\parallel\mathrm{LP}(d_{\rm ovk})\parallel\mathrm{LP}(M),\\
 \sigma&\leftarrow\mathcal S.\mathsf{Sign}(\mathsf{sk}_s,\operatorname{EncMsg}(r,E)),\\
 P&=\mathsf D_{\rm sig}\parallel\mathrm{LP}(\sigma)\parallel\mathrm{LP}(E).
\end{align*}
混合加密得到内层密文字节$B$后，用一次性私钥对上下文、公钥和$B$签名：
\[
 \tau\leftarrow\mathsf{OTS.Sign}_{\mathsf{osk}}
    (\mathsf D_{\rm ct}\parallel\mathrm{LP}(\mathsf{ctx})
      \parallel\mathrm{LP}(\mathsf{ovk})\parallel\mathrm{LP}(B)).
\]
最终传输的是$(\mathsf{ctx},\mathsf{ovk},B,\tau)$的规范编码。接收方先检查这层签名，再解密$B$；恢复出的IBS签名还必须认证同一个公钥承诺。这样，外层字节认证与内部身份认证建立了联系。一次性私钥仅用于当前密文，使用后释放。

\begin{figure}[H]
\centering
\begin{tikzpicture}[node distance=8mm and 8mm]
 \node[flow,text width=42mm] (input) {消息$M$＋接收身份$r$\\一次性公钥承诺};
 \node[flow,right=of input,text width=42mm] (sig) {发送者IBS签名\\形成载荷$P$};
 \node[flow,below=of input,text width=42mm] (seed) {新32字节种子$z$\\HIBE按$(r,s)$加密};
 \node[flow,below=of sig,text width=42mm] (dem) {由$z$派生密钥\\GCM加密载荷$P$};
 \node[check,below=of dem,text width=42mm] (outer) {对全部内层字节签名\\输出一次性认证外帧};
 \draw[flowarrow] (input)--(sig);
 \draw[flowarrow] (sig)--(dem);
 \draw[flowarrow] (seed)--node[above,font=\small]{$K$}(dem);
 \draw[flowarrow] (dem)--(outer);
 \draw[flowarrow] (seed.south)|-(outer.west);
\end{tikzpicture}
\caption{从身份与应用消息生成完整密文；HIBE胶囊和对称密文一起受到外层认证}
\label{fig:protocol-generation}
\end{figure}

\section{接收端的完整处理顺序}
接收方按从外到内的顺序处理密文：
\begin{enumerate}
 \item 检查长度、版本、字段数量和公开上下文，验证外层一次性签名。
 \item 加载内层密文，以接收密钥$r'$为期望发送者$s'$委派叶钥。
 \item 用叶钥解开HIBE胶囊，恢复32字节种子并导出对称密钥。
 \item 验证GCM标签，解析内部IBS签名、一次性公钥承诺与消息。
 \item 以$s'$验证绑定$r'$的IBS消息，再核对内外一次性公钥承诺。
 \item 全部检查通过后输出$M$；数据错误统一返回拒绝结果$\bot$。
\end{enumerate}
解析采用明确的域标识和长度前缀；截断、追加、重复字段、非规范整数和错误对象类型均被拒绝。依赖缺失、资源超时等执行故障单独提示，便于使用者处理。

\section{三组实例与所用原语}
表~\ref{tab:representatives}列出报告重点分析的三组组合。三组共享上述协议，差别在于HIBE、IBS及对应参数。应用容量同时受IBS签名输入和混合载荷预算限制。

\begin{table}[H]
\centering\small
\caption{三组代表及应用容量}
\label{tab:representatives}
\begin{tabularx}{\textwidth}{@{}cp{.16\textwidth}Xrr@{}}
\toprule
\rowcolor{black!8} 组 & HIBE & IBS & 消息上限/B & 签名预算/B \\
\midrule
A & LATTE-3 & BNN/ML-DSA-44 & 59,191 & 6,237 \\
B & LATTE-3 & SRGG PARAM-I紧凑格式 & 65,364 & 15,704 \\
C & Lai C2 & Sageloli PARAM-I & 3,924 & 3,227,824 \\
\bottomrule
\end{tabularx}
\end{table}

A、B的混合载荷预算分别为64 KiB和128 KiB。C的签名约3.08 MiB，采用8 MiB预算；其3,924字节应用上限由IBS的4,096字节输入限制扣除协议编码得到。三组的HIBE输入均为32字节种子。

\subsection{LATTE：接收身份与发送身份进入两层陷门}
LATTE在负循环环$R_q=\mathbb Z_q[x]/(x^N+1)$上建立陷门~\cite{zhao2024latte}。本作品采用$N=1024$、$q=68718428161$的LATTE-3参数。建立阶段取短多项式$f,g$，求解
\[
 fG-gF=q,\qquad h=gf^{-1}\pmod q,
 \qquad S_0=\begin{pmatrix}g&-f\\G&-F\end{pmatrix}.
\]
$S_0$是根陷门。接收身份经域分隔哈希生成$A_1=H_1(r)$，根陷门委派为一级基；发送身份进入$A'_2=H_2(r,s)$，由一级基提取叶钥$(t_0,t_1,t_2)$。叶钥满足
\[
 t_0+t_1h+t_2A_1=A'_2\pmod q.
\]
该关系把叶钥与完整身份路径联系起来，也是加载密钥时的检查依据。实现使用作者的Gaussian采样和ffSampling路径，修订COSAC读取新Box--Muller缓冲的位置，并对原生输出执行精确环关系检查。HIBE解密保留恢复种子后的重加密核对。

\subsection{Lai：追加身份块完成两级委派}
Lai方案以理想格陷门组织委派~\cite{lai2015trapdoors}。公开向量与短陷门满足
\[
 \mathbf a_0^{\mathsf T}R+\mathbf a_1=\mathbf g.
\]
两层身份分别映射为二进制环元素$\mathrm{id}_1,\mathrm{id}_2$，形成
\[
 \mathbf a_{\mathrm{id}_1}=\mathbf u_0+\mathrm{id}_1\mathbf u_1,
 \qquad \mathbf a_{\mathrm{ID}}=\mathbf a_{\mathrm{id}_1}+\mathrm{id}_2\mathbf u_2.
\]
每层向公开矩阵追加对应身份块，并生成新的短陷门。加密取短矩阵$R_0,R_1,R_2$，令
\[
 R_{\mathrm{ID}}=R_0+\mathrm{id}_1R_1+\mathrm{id}_2R_2,
 \qquad\mathbf e_1=-\mathbf e_0^{\mathsf T}R_{\mathrm{ID}},
\]
再生成
\[
 \mathbf u=F_1\mathbf s+p\mathbf e_0,\qquad
 \mathbf v=(\mathbf a_{\mathrm{ID}}+h'\mathbf g)\mathbf s+p\mathbf e_1+\mathbf m.
\]
解密使用叶陷门完成Babai反演，检查误差的模$p$结构并恢复消息。代表C采用$n=2048$、$q=274877976577$、$k=39$、深度2、$p=2$的C2配置，噪声参数$\sigma=2$、$\sigma_r=1$。实现按层追加身份块，依据实际方程统一矩阵形状，陪集采样使用Klein/GPV路径。身份映射取二进制系数，以保持身份乘短矩阵后的噪声大小。

\subsection{BNN/ML-DSA：机构证书与用户签名共同认证身份}
BNN转换由机构为身份$s$生成用户签名公钥$\mathsf{vk}_s$，并签发证书~\cite{galindo2006ibs}：
\[
 \mathsf{cert}_s=\mathsf{DS.Sign}_{\mathsf{sk}_{\rm CA}}
     (\mathsf{Enc}(s,\mathsf{vk}_s)).
\]
用户再用私钥签署包含上下文、身份、公钥、证书和协议消息$M'$的编码；IBS签名携带用户公钥、证书和用户签名。验证时先检查证书，再验证用户消息。代表A采用ML-DSA-44~\cite{nist2024fips204}，其基础公钥为1,312字节，基础签名为2,420字节；组合编码后的IBS签名为6,237字节。

\subsection{SRGG：身份目标与消息挑战进入同一验签关系}
SRGG采用模块格结构，公开矩阵为$A=[I_d\mid\widehat A\mid B]$~\cite{jia2025ibs}。身份决定目标$h=H_1(P,s)$，用户短见证$y$满足$Ay=h\pmod Q$。签名构造
\[
 M_h=[\,2A\mid Qj-2h\,]\pmod{2Q},\qquad
 z=p_0+\begin{bmatrix}y\\1\end{bmatrix}(\zeta u+c),
\]
其中$\zeta=1+x^{n/2}$、$j=\zeta^\star e_1$，扰动$p_0$的承诺和消息共同决定挑战$c$。验证者由$W'=M_hz-Qjc$恢复承诺，重新计算挑战，并检查响应范数。代表B采用$n=512$、模块秩2、$p=16383$、$Q=15744063$、gadget基31且位数2的PARAM-I。第3章介绍其对象编码。

\subsection{Sageloli：用环陷门生成身份签名}
Sageloli的公开向量与主陷门$R$满足~\cite{sageloli2023ibs}
\[
 A=[\bar A\mid g-\bar AR].
\]
发钥使用身份与32字节盐$\rho$计算$h=H_1(P,\rho,s)$，并采样身份陷门$T_s$，满足$AT_s=h-g\pmod q$。签名另取消息盐$\eta$，通过相关扰动和三进制陪集采样产生响应$z$。验证者检查范数和
\[
 [A,h]z=H_2(P,\rho,s,\eta,M')\pmod q.
\]
代表C采用PARAM-I：$(d,k,\ell)=(2048,65,132)$、$q=3^{65}$、$r^2=49$；身份与签名采样宽度平方分别为$2720281532542$和$270116709797628319132317$。第3章重点说明怎样在该规模执行采样。

\section{参数与安全依据}
作品采用Wang通用构造的普通隐私、加强认证与自我伪造防护目标，两层HIBE采用身份匿名性条件；具体困难假设、参数和哈希模型以相应原文为依据。哈希接口以SHAKE等实现，按论文的随机预言机模型说明。三代表采用上述论文或标准参数，调试参数不进入本报告实验。

ML-DSA-44的底层标准参数对应NIST安全类别2；SRGG PARAM-I论文标称目标为100 bit，并给出量子密钥恢复与伪造成本95 bit和93 bit。Lai C2的系数实例估计记录为经典约169.9 bit、量子约154.2 bit；Sageloli采用原文PARAM-I。各数值对应不同部件和模型，因此第5章以完整实现成本解释各组特点。功能测试用于核对实现，不能代替安全性论证。
